\chapter{Eye Discharge}

\section*{Definition}
Abnormal secretion from the eye ranging from clear and watery to thick, purulent, or mucoid, indicating underlying ocular surface or lacrimal system pathology.

\section*{Pathophysiology}
Watery discharge typically accompanies viral conjunctivitis or allergic reaction with increased lacrimation. Purulent (yellow/green) discharge indicates bacterial infection with neutrophil accumulation. Mucoid discharge (stringy, white) suggests allergic conjunctivitis. Chronic discharge may indicate lacrimal duct obstruction or dacryocystitis.

\section*{Etiology}
\begin{itemize}
  \item Conjunctivitis (viral, bacterial, or allergic)
  \item Dacryocystitis (lacrimal sac infection)
  \item Nasolacrimal duct obstruction
  \item Keratitis (corneal infection or inflammation)
  \item Stye (hordeolum) or chalazion
  \item Blepharitis (eyelid margin inflammation)
  \item Entropion or ectropion (eyelid malposition)
  \item Foreign body
  \item Contact lens-related infection
  \item Herpes simplex or herpes zoster ophthalmicus
\end{itemize}

\section*{Indications for Evaluation}
Seek care if:
\begin{itemize}
  \item Severe or worsening symptoms
  \item Eye discharge with vision changes, severe eye pain, photophobia, or in contact lens wearers (risk of corneal ulcer)
  \item Accompanied by other concerning signs
\end{itemize}

\section*{Related Symptoms}
\begin{itemize}
  \item Eye redness
  \item Itchy eyes
  \item Blurred vision
  \item Eye pain
  \item Photophobia
\end{itemize}
\textbf{Note:} Always consider serious underlying conditions when evaluating persistent eye discharge.

\section*{Self-Care / Home Management}
\begin{itemize}
  \item Rest and avoid activities that may aggravate the condition.
  \item Maintain adequate hydration and a balanced diet.
  \item Over-the-counter remedies may provide symptomatic relief (consult a pharmacist or doctor).
  \item Monitor symptoms and seek professional advice if they persist or worsen.
  \item Avoid self-medicating with prescription medications without medical guidance.
\end{itemize}

\section*{Diagnostic Approach}
\begin{itemize}
  \item A thorough history and physical examination are the first steps in evaluation.
  \item Laboratory tests (blood work, urinalysis) may be ordered based on clinical suspicion.
  \item Imaging studies (X-ray, ultrasound, CT, MRI) may be indicated when structural pathology is suspected.
  \item Specialist referral is arranged when the diagnosis remains uncertain or requires specific expertise.
\end{itemize}

\section*{Treatment / Management Overview}
\begin{itemize}
  \item Treatment targets the underlying cause identified through diagnostic evaluation.
  \item Supportive care and symptom management are provided while awaiting definitive therapy.
  \item Medications, lifestyle modifications, and physical therapies may be prescribed as appropriate.
  \item Follow-up ensures treatment effectiveness and allows timely adjustments to the care plan.
\end{itemize}

\clearpage